% Part: applied-modal-logics
% Chapter: temporal-logic
% Section: introduction

\documentclass[../../../include/open-logic-section]{subfiles}

\begin{document}

\olfileid[zh]{aml}{tl}{int}

\olsection{导言}

时态逻辑处理关于将会是实情、抑或曾经是实情之事物的断言。\zhFullName{Arthur Prior}{亚瑟·普莱尔}{普莱尔}被认为是时态逻辑的创始者，他称之为「时态逻辑」（\textit{tense logic}）。这里我们对时态逻辑的处理，大体上遵循\zhFullName{Arthur Prior}{亚瑟·普莱尔}{普莱尔}原先的模态处理方式，即在命题逻辑的基本框架中引入时态算子；命题逻辑把断言一般视为缺乏时态的。

例如，在命题逻辑中，我或许会谈论一条叫 Beezie 的狗，它有时坐着、有时不坐，狗往往就是这样。在经典逻辑中，断言 Beezie 坐着同时又断言 Beezie 不坐着会是矛盾的。但显然，这两者可以同时为真，只是不能在同一时刻为真；在语言中加入时态算子，就能让我们相当容易地表达这一断言。时态算子的加入还使我们得以说明某些推断的有效性，例如从「Beezie 将得到一份零食或一个球」推出「Beezie 将得到一份零食，或 Beezie 将得到一个球」。

然而，时态逻辑会引发许多哲学问题，这些哲学问题或许会促使我们在不同时态逻辑框架之间作出取舍。例如，一个未来偶然命题是关于未来的陈述，它既非必然亦非不可能。如果我们说「Richard 明天会去食品杂货店」，那么我们所表达的乃是关于一件尚未发生之事的断言，而该断言的真值是可争议的。事实上，这一断言究竟能否首先被\emph{赋予}真值，本身就是可争议的。如果我们持严格的决定论立场，那么我们或许可以坦然接受如下观念：即便在所涉事件被认为发生之前，这一语句事实上也已为真或为假——只是我们或许尚不知道它的真值。相比之下，我们或许相信一个真正开放的未来，在其中未来偶然命题的真值是未定的。

事实证明，这些关于时间的结构与本性的诸多立场，都被内嵌进我们对时态逻辑之模型与框架的选择之中。例如，我们或许会问自己：是否应当构造时间呈线性、分叉乃至环形的模型？我们或许不得不就时态模型是否有起点与终点作出决定，以及时间究竟该用离散的瞬间还是作为一种连续统来表征。

\end{document}
